% Generated draft: requires mathematical and source review.
\documentclass[11pt]{article}
\usepackage[margin=25mm]{geometry}
\usepackage{fontspec}
\setmainfont{Latin Modern Roman}
\usepackage{amsmath,amssymb,mathtools,unicode-math}
\setmathfont{Latin Modern Math}
\usepackage{enumitem,needspace,xurl,hyperref}
\setlength{\parindent}{0pt}
\setlength{\parskip}{5pt}
\setlength{\emergencystretch}{3em}
\newcommand{\N}{\mathbb{N}}
\newcommand{\Z}{\mathbb{Z}}
\newcommand{\Q}{\mathbb{Q}}
\newcommand{\R}{\mathbb{R}}
\newcommand{\C}{\mathbb{C}}
\newcommand{\F}{\mathbb{F}}
\newcommand{\A}{\mathbb{A}}
\newcommand{\PP}{\mathbb{P}}
\newcommand{\E}{\mathbb{E}}
\newcommand{\eps}{\varepsilon}
\newcommand{\dd}{\,\mathrm{d}}
\begin{document}
\section*{UnsolvedMath ID: 30006496}
\textbf{Source problem code:} catalogue entry
\subsection{Definitions and mathematical statement}
We use the finite-dimensional complex operator interpretation of the supplied statement. The record does not explicitly specify the scalar field or require finite-dimensionality, so these conventions require confirmation against the source before acceptance.

Let \(V\) be a nonzero finite-dimensional complex Hilbert space. A unitary Yang--Baxter operator on \(V\) is a linear map \(R:V\otimes V\to V\otimes V\) satisfying
\[
R^*R=RR^*=I_{V\otimes V}
\]
and the braid-form Yang--Baxter equation
\[
(R\otimes I_V)(I_V\otimes R)(R\otimes I_V)
=(I_V\otimes R)(R\otimes I_V)(I_V\otimes R)
\quad\text{on }V^{\otimes 3}.
\]
For each integer \(n\geq 2\), let \(B_n\) be the braid group with generators \(\sigma_1,\ldots,\sigma_{n-1}\) and relations
\[
\sigma_i\sigma_{i+1}\sigma_i=\sigma_{i+1}\sigma_i\sigma_{i+1}
\quad(1\leq i\leq n-2),
\qquad
\sigma_i\sigma_j=\sigma_j\sigma_i
\quad(|i-j|\geq 2).
\]
The associated representation is
\[
\rho_{R,n}:B_n\longrightarrow U(V^{\otimes n}),
\qquad
\rho_{R,n}(\sigma_i)
=I_V^{\otimes(i-1)}\otimes R\otimes I_V^{\otimes(n-i-1)},
\]
where a zeroth tensor power of an identity is omitted. Write \(G_{R,n}=\rho_{R,n}(B_n)\). Following the source's explicit formulation, virtual abelianity means that there exists an abelian normal subgroup \(A\trianglelefteq G_{R,n}\) with finite group index \([G_{R,n}:A]\).

\textbf{Supplied decategorified conjecture, under the conventions above.} For every such \(V\), every unitary Yang--Baxter operator \(R\) on \(V\), and every integer \(n\geq 2\), does there exist an abelian normal subgroup \(A_{R,n}\trianglelefteq G_{R,n}\) such that
\[
[G_{R,n}:A_{R,n}]<\infty?
\]
The subgroup and its index may depend on \(V\), \(R\), and \(n\). The image here is the abstract group image, without taking a topological closure. If \(n=1\) is included, \(B_1\) and its image are trivial. Whether infinite-dimensional Hilbert spaces are intended is unresolved in the supplied record; the displayed formulation does not address that extension.
\subsection{Short English statement}
The supplied conjecture asks whether every braid-group representation obtained by placing a fixed unitary Yang--Baxter operator on adjacent tensor factors has an abelian normal subgroup of finite index in its image, for every number of strands. The finite-dimensional complex Hilbert-space convention needs source confirmation. The record labels the general conjecture open; that label is supplied dated curation, not an independently verified assessment.
\subsection{Sources}
\begin{itemize}
\item Supplied local UnsolvedMath registry record, permanent ID \(30006496\), source problem code \texttt{catalogue entry}, title: \emph{Virtual Abelianity of Unitary Yang--Baxter Braid Images}. The project identifies \texttt{data/unsolvedmath/problems.json} and \texttt{catalog/problems.json} as local registries; this draft uses the supplied complete record without independently inspecting those files.
\item Supplied report citation: \emph{Mini-Workshop: The Yang--Baxter Equation and Representations of Braid Groups (2026)}, Oberwolfach Reports, pp.\ 2633--2672, DOI \texttt{10.4171/owr/2025/49}. \url{https://doi.org/10.4171/owr/2025/49}. The supplied excerpt bears page number 2665 and calls the statement a ``Decategorified Conjecture,'' referring to its internal reference [3]. No proposer or bibliographic details for [3] are supplied. The citation's 2026 year differs from the report identifier and proposed year 2025.
\item Supplied literature source: \emph{Oberwolfach Report 49/2025}. \url{https://doi.org/10.4171/owr/2025/49}.
\item Supplied literature source: \emph{Localization of unitary braid group representations}. \url{https://arxiv.org/abs/1009.0241}.
\item Literature assessment dated exactly \texttt{2026-08-22}: the supplied curation labels the general conjecture open, describes substantial evidence and special-case proofs, and says the general statement remains unresolved. This is supplied dated curation, not an independently verified literature assessment. The background explicitly requests independent verification of time-sensitive status and cited claims.
\item The registry marks the statement \texttt{corrected\_verified} and reports that the official report was checked and notation repaired. These are supplied provenance claims; no source visit or independent verification is asserted here.
\end{itemize}
\end{document}
